\hypertarget{chapter-13-the-sync-that-lied}{%
\section{Chapter 13: The Sync That
Lied}\label{chapter-13-the-sync-that-lied}}

\begin{quote}
\itshape ``The cruelest lie is not \emph{you are wrong}. It is \emph{you remember
too much}.''\\
\normalfont\hfill---verify-before-sync field note
\end{quote}

\hypertarget{scene-1-the-shorter-history}{%
\subsection{Scene 1: The Shorter
History}\label{scene-1-the-shorter-history}}

Of all the ways to kill a thing that cannot be bribed or counterfeited, the
oldest and quietest is to convince it that less happened than it remembers.

Elena Voss had seen people destroyed this way---assets persuaded that the years
they had given had never been logged, operatives told the missions they
remembered were administrative errors. Memory was the softest target in any
system, because a system that doubts its own past will accept almost any present
offered to replace it.

The attack on the Sigil came dressed as good news. A peer appeared on the
network---well-connected, fast, generous---announcing a chain. Not a longer one.
A \emph{shorter} one. It claimed, with perfect cryptographic politeness, that the
true history was thousands of blocks behind where every honest node believed it to
be, and it offered, helpfully, to bring everyone back into agreement.

``If a node believes it,'' Wren said, very quietly, watching the announcement
propagate, ``it doesn't just adopt a wrong number. It \emph{deletes}. It throws
away every block past the lie, and every balance those blocks recorded, and calls
it housekeeping. Thousands of blocks of real history, gone, because a stranger
said \emph{you remember too much}.''

Elena understood the horror of it immediately, because it was the exact shape of
the thing she had spent her life resisting. ``It's not stealing the money,'' she
said. ``It's stealing the \emph{proof the money was ever yours}.''

\hypertarget{scene-2-the-handshake}{%
\subsection{Scene 2: The
Handshake}\label{scene-2-the-handshake}}

The defense, Wren explained, was a rule so simple it sounded like cowardice, and
so important that the whole network's survival rested on it: \emph{a node must
never, under any circumstances, follow a peer downward.}

``Forward, with proof, gladly,'' she said. ``If a peer has more history than I do,
I demand it show me---block by verifiable block, each one valid, each one building
on the last---and only then do I adopt it. But \emph{backward}? Never. A peer that
asks me to discard history I can prove I hold is, by definition, either broken or
hostile. There is no honest reason to ask someone to remember less. The node
doesn't argue. It simply refuses to shrink.''

It was a handshake before a sync---a verification before trust, the same liturgy
as everything else here. \emph{Show me yours is longer and every block of it is
true, or I keep mine.} The poisoned peer's shorter chain, however politely
offered, could never satisfy it, because shorter was the one thing the rule would
not accept.

``So why is anyone worried?'' Elena asked. ``If the rule holds, the lie just
bounces off.''

``Because rules live in code,'' Wren said, ``and someone is about to argue that
this one is too strict.''

\hypertarget{scene-3-the-reasonable-exception}{%
\subsection{Scene 3: The Reasonable
Exception}\label{scene-3-the-reasonable-exception}}

The release came through the quorum at the worst possible moment, when half the
network was distracted defending against the false peer. It was, of course,
beautifully reasonable. It proposed a small relaxation: that in cases of
``severe network partition,'' to ``aid recovery,'' a node might be permitted to
follow a peer to a \emph{slightly} lower height, if enough peers agreed on it.

A tolerance. Again. The same blade, the same handle, worn smooth now by Elena's
recognition.

``They're not attacking the rule,'' she said. ``They're attacking the word
\emph{never}. Turn \emph{never} into \emph{usually}, and the false peer doesn't
need to convince a node anymore. It just needs to bring enough friends to vote
that everyone remembers too much.''

She did what she had learned to do. She did not plead. She published a
proof---plain, reproducible, gossiped to the whole network---demonstrating
exactly what the relaxation permitted: a path, narrow but real, by which a
coordinated set of liars could march the entire chain backward and erase whatever
history stood between the truth and their preferred past. Not a recovery
mechanism. A demolition permit with good manners.

The network read the proof, and the network did the arithmetic, and the
relaxation died at five signatures with a sixth recoiling visibly. \emph{Never}
survived, as a word and as a wall.

\hypertarget{scene-4-what-a-chain-is-for}{%
\subsection{Scene 4: What a Chain Is
For}\label{scene-4-what-a-chain-is-for}}

The false peer, finding no node willing to shrink and no rule willing to bend, did
what lies do when no one believes them. It withered. Its announcements grew
quieter, then stopped, the way a rumor stops when the last person who might have
repeated it decides to check first.

Later, on the rooftop that had become their place of after, Wren said: ``You know
what you actually defended tonight. It wasn't the balances.''

``It was the \emph{continuity},'' Elena said. She was looking east, where the sky
was the gray of a screen waking up. ``The thing the Veil never had. A history that
can grow but cannot be made to forget. That's what a chain is \emph{for}---not to
remember everything, anyone can remember. To make forgetting require everyone's
honest consent, and to make dishonest forgetting impossible to hide.''

She thought of the burned years of her old life, the missions unlogged, the proof
erased by people who owned the records. She had spent two years certain that the
best a person could do was carry the truth alone and hope to outlive the people
who denied it.

She had been wrong about that too. You did not have to carry it alone.

You could put it somewhere that refused to shrink.

\hypertarget{chapter-13-notes}{%
\subsection{Chapter Notes}\label{chapter-13-notes}}

\begin{itemize}
\tightlist
\item
  \textbf{The sync-down attack.} The chapter dramatizes a real catastrophic
  failure mode: a node persuaded that the canonical chain is \emph{shorter} than
  its own may discard valid history and the balances it recorded. The danger is
  data loss, not theft of a key.
\item
  \textbf{Never follow a peer downward.} The defense is an asymmetry: adopt a
  longer chain only after verifying every block, but \emph{never} shrink to match
  a shorter one. There is no honest reason to ask a node to remember less.
\item
  \textbf{Verify before you sync.} A handshake that demands proof of a longer,
  valid history before adoption is the same ``verify before you trust'' liturgy
  applied to history itself.
\item
  \textbf{The word ``never'' as the real target.} The recurring antagonist
  doesn't attack the rule head-on; it proposes a reasonable-sounding exception
  that converts ``never'' into ``usually,'' which is all a coordinated attacker
  needs. Defending the absoluteness of the rule is the chapter's quiet climax.
\end{itemize}
